\section{Proposed MoM Bundle Format}
\label{app:artifact-profile}
\begingroup
\footnotesize

This appendix instantiates the artifact contract in
Section~\ref{sec:artifact}. The format is a versioned proposal; a
machine-readable companion supplies schemas for the manifest, typed IR,
realization profile, binding, lock, evaluation protocol, evaluation
attestation, and run record, together with one verified example.

\subsection{Bundle contents and evidence boundary}

The manifest is an immutable index over the human model coordinate, behavior
variant, semantic digest, typed entrypoints, canonical specification, semantic
assets, contracts, evaluation protocols, and optional realization templates.
Canonical IR remains the sole semantic authority. Authoring source and model
cards are optional distribution objects rather than executable authorities.
Evaluation results, signatures, SBOMs, and build provenance are separate
digest-bound attestations: publishing new evidence does not change the model
being evaluated.

Every object descriptor contains a normalized relative path, media type,
SHA-256 digest, and byte size. Absolute, backslash, empty, dot, repeated,
traversal, and symlink paths are rejected before parsing. A production archive
importer must additionally bound expanded size, file count, nesting, and
compression ratio and reject hard links, devices, and other non-regular files.

The example exposes one accuracy-first model and two realization profiles:
ROCm plus cloud and CUDA plus cloud. Behavior variants belong to logical
identity; realization profiles contain endpoint-free compatibility
requirements and do not. A binding chooses one realization profile and supplies
concrete candidates.

\subsection{Four identities without circular hashes}

Let $C$ be canonical JSON and $H$ be SHA-256. Let $I$ be the fully defaulted and
typed IR, whose semantic assets are themselves content-addressed. Let $M$ be
the manifest, whose object descriptors commit to the distributed byte closure.
Let $N(B_s)$ be the normalized typed representation obtained from binding
source $B_s$, and let $L$ be the completed normalized lock. The four identities
are
\begin{equation}
 \begin{aligned}
 d_{\mathrm{semantic}} &= H(C(I)),&
 d_{\mathrm{bundle}}   &= H(C(M)),\\
 d_{\mathrm{binding}}  &= H(C(N(B_s))),&
 d_{\mathrm{lock}}     &= H(C(L)).
 \end{aligned}
 \label{eq:bundle-identities}
\end{equation}
No object stores its own digest. The manifest stores
$d_{\mathrm{semantic}}$; the binding stores semantic and bundle subjects; the
lock stores that subject and $d_{\mathrm{binding}}$; traces and attestations
then store all four. Formatting a YAML binding therefore cannot create a new
environment identity.

\subsection{Target compatibility and locking}

A logical constituent declares an identity and required or preferred
capabilities, including protocol, modality, interface features, input and
output limits, and data policy. A realization profile may additionally require
a wire transport, API adapter, serving runtime, leaf-artifact format,
accelerator API, vendor, architecture, or deployment driver. Environment policy may strengthen but
never relax semantic constraints. Resolution is thus the intersection of
logical admissibility, realization-profile requirements, and environment
policy.

The lock records the complete candidate inventory available to an execution,
not the one candidate selected by a particular request. For open artifacts it
may pin a content digest or repository revision together with tokenizer,
runtime image, protocol adapter, and accelerator evidence. For a closed service
it marks identity as observed-opaque, records an observation time and probe
digest, and makes no bitwise-reproducibility claim. Floating dependencies are
excluded from a production lock.

\subsection{Distribution closures and evaluation}

A thin export retains immutable external references; a materialized export
vendors leaf artifacts whose licenses permit redistribution; and a targeted
export additionally carries runtime objects for selected realization
profiles. Materialization fails for a closed provider rather than treating an
API alias as an offline model. A canonical directory and deterministic archive
define local exchange; OCI descriptors, manifests, and referrers provide a
natural registry transport and evidence attachment.

An evaluation protocol is part of the bundle because it defines an assessment
procedure. A result is an external attestation over semantic, bundle, binding,
lock, protocol, dataset, and harness digests, plus sampling and concurrency
settings. Serving and evaluation therefore consume the same resolved model
identity, while later measurements remain independently publishable.

\subsection{Conformance}

The companion verifier checks all supplied schemas; stable canonical
identities; traversal-safe paths, byte sizes, and digests; manifest/IR semantic
closure; behavior variants and realization profiles; graph reachability,
explicit contract branches, universal termination, and finite bounds; fail-closed hard
constraints; constituent identity and capability satisfaction; complete lock
inventory; pinned and observed-opaque evidence; secret non-inclusion; and run
and evaluation attribution to all four identities.

A complete cross-engine suite must additionally test deterministic compilation,
source-to-IR semantic preservation, operator capability negotiation, archive
hardening, license and extension trust, trace conformance, and portability
across at least two disclosed bindings while fixing semantic and bundle
digests.
\endgroup
